\section{My impact}

\subsection*{3a. How my contributions drove the team forward}

The honest answer is that my simulation-first instinct in wk18 is what kept the project on the rails. I want to defend that claim with evidence rather than assertion, because it is exactly the kind of thing that at Merit-level would read as self-flattery.

The criterion I used to judge my own effectiveness is the simplest one available: against the 12 requirements in the brief (R01--R12), did the team have a working path to each by the submission deadline? Against that measure, by wk22 we had 12/12 implemented in simulation and verified end-to-end in \texttt{simple\_simulator.py}, and we had 4/12 implemented on the Pi bench (limited by hardware access, not code). The 12/12 in simulation is what matters here. R01 (flight-area geofence), R02 (SSSI no-overfly), R03 (TOL distance check), R04 (50~m altitude cap), R05 (search pattern + detection), R06 (PLB-triggered focus area), R07 (land 5--10~m from casualty), R09 (RTH + motor cutoff), R10 (detection photos + JSON), and R12 (MIT-licensed public repo) were all closed in code that I wrote or drove. The evidence is the commit graph: 820+ commits on my branches, and the single audit commit \texttt{ce5c036} that walked the brief requirements one by one and closed the gaps I found. That audit alone caught R03 and R07 as silent gaps --- the code technically arm-disarmed fine but never checked distance from TOL before arming, and it logged landing distance from home rather than from the casualty. These are the kind of small, correct-sounding defaults that fail quietly in a client demo.

Where I drove the team forward most visibly, though, was in unblocking others. Robin's flight script only works because I built him the passive-watch HTTP integration (\texttt{a5b1ab7}), wrote him the parsing examples (\texttt{885cee4}), and added the \texttt{--smart-clear} flag he needed to re-lock after the first target (\texttt{cec899e}). Demetro's Final Design Review presentation only works because I built his two slides from scratch (\texttt{32fb0c5}) and then wrote his speaking scripts (\texttt{b2a9cd4}, \texttt{ba667f2}) --- 465 words, 3.1 minutes of tested delivery --- because he asked for help framing the technical material for a non-specialist audience. Neither of those was glamorous work. Both of them were the difference between his deliverable shipping on time and shipping late.

The forward-looking part of this honest evaluation is harder. My measurable individual velocity was high --- 820+ commits, 2508 lines of simulator, 900 lines of main.py, 127 unit tests --- but my team multiplier was uneven. I was a force multiplier for Robin and Demetro directly (via integration work and slide-building) and a zero multiplier for Edward and [PLACEHOLDER: hardware-lead], because I never sat down with them on the hardware side. I evaluated my own impact against a second, harder criterion: if I had been hit by a bus in wk18, what would have survived? The \texttt{simple\_simulator.py} would have survived as a working MVP, the retrained model would have deployed cleanly (commit message says ``drop-in replacement, no code changes''), and the test ladder would have remained a readable list of scripts to run. The state machine and the geofence math would \emph{not} have survived, because they lived almost entirely in my head and in an informal line-range blueprint. That is a brittleness I created by soloing too hard in wk14--wk17, and it is the specific thing I would unwind in a future project.

\subsection*{3b. What I would add or re-focus on in future}

Three shifts, each tied to a specific observation from this project.

First, I would \textbf{teach, not ship}. The largest gap between my individual output and my team impact was that I shipped a lot of working code that nobody else read. In a five-person team, the multiplier on a 100-line function you walk a teammate through is higher than the multiplier on a 500-line module you commit alone. I now see that the engineering I am best at --- building self-contained tools fast --- is also the engineering that has the lowest team multiplier, and that I was using my strength as an excuse to avoid the thing I am worse at, which is pair programming with people whose style differs from mine. The mechanism I am committing to for the next project is: no module larger than 300 lines without at least one live code-walkthrough with a teammate before merge. The observable signal I will watch for to know it is working is whether any teammate other than me edits one of my files in a given week. In this project that number was zero for \texttt{simple\_simulator.py} and \texttt{main.py}, which is the evidence I need to take the lesson seriously.

Second, I would \textbf{scope for the team, not for the problem}. My FSM in \texttt{main.py} has 32 state transitions, which is appropriate for an autonomous SAR mission but inappropriate for a five-person student team with limited shared context. In hindsight I over-engineered the state machine in exactly the way Robin called out in the wk17 disagreement, and I did it because I was thinking about the problem (``what does the drone need to do?'') and not about the organism (``what can this team maintain?''). The Hatton \& Smith ``critical reflection'' move here is to notice that my bias toward full-coverage design is partly a risk-aversion move on my own part --- I build comprehensive things because I am uncomfortable leaving edge cases unhandled, and that discomfort is stronger than my awareness of the maintenance cost I am imposing on others. The mechanism I am committing to is: before committing any state machine with more than 10 states, I will present the state diagram to a teammate and let them cut it. If they cannot cut anything, the diagram is fine. If they cut something, I have to defend the version I had, or adopt theirs.

Third, I would \textbf{build feedback channels for my own work earlier}. The thing that most improved the quality of my output in this project was the D7 scoring-rubric pass I applied to my own reflective drafts (six iterations, 68 $\rightarrow$ 83--84) and the requirements-matrix audit I ran on the codebase in wk22 (commit \texttt{ce5c036}). Both of these were structured external evaluations I set up for myself because I did not trust my own judgement in isolation, and both of them caught things I would not have caught alone. The generalised lesson is that solo work scales until you hit the limit of your own self-assessment, at which point you need either another human or a structured scoring rubric to keep improving. I hit that limit around wk20 on the vision code and did not notice until the geofence RTL crash in wk22 exposed it. I now see that the right time to put a rubric on your own work is \emph{before} the crash, not after it.
